\documentclass{article}
\usepackage[a4paper,margin=2.5cm]{geometry}
\usepackage{amsmath, amssymb, amsthm}
\usepackage{booktabs}
\usepackage{graphicx}

\newtheorem{proposition}{Proposition}
\newtheorem{lemma}[proposition]{Lemma}
\newtheorem{corollary}[proposition]{Corollary}
\theoremstyle{remark}
\newtheorem*{remark}{Remark}

\title{Taylor expansion of the optimal winning probability $w_{n,p}$ about $p=\tfrac12$}
\author{}
\date{2026-04-16}

\begin{document}
\maketitle

\section{Setting}

We consider the all-heads-wins game: start with $n$ coins, each landing
heads independently with probability $p\in(0,1)$; in every round all
remaining coins are tossed and at least one must be set aside (never
tossed again); the game is won iff every set-aside coin shows heads. Let
$w_{n,p}$ be the winning probability under optimal play, with
$w_{0,p}:=1$. The Bellman recursion is
\begin{equation}\label{eq:bellman}
  w_{n,p}\;=\;p^{n}+\sum_{j=1}^{n-1}\binom{n}{j}\,p^{n-j}(1-p)^{j}\,
  M_{n,j}(p),
  \qquad M_{n,j}(p)\;:=\;\max_{j\le m\le n-1}w_{m,p}.
\end{equation}
We write $w_n$ for $w_{n,p}$ when $p$ is understood.

At $p=\tfrac12$ every strategy gives value $\tfrac12$, hence
$w_{n,\,1/2}=\tfrac12$ for all $n\ge 1$.

\paragraph{Caveat on regularity.}
The function $p\mapsto w_{n,p}$ is continuous but typically \emph{not}
differentiable at $p=\tfrac12$ once $n\ge3$: the $\max$ in
\eqref{eq:bellman} is a two-sided maximum of polynomial pieces, so there
is a corner at $p=\tfrac12$ whose left- and right-derivatives disagree in
general. We therefore introduce \emph{one-sided} Taylor expansions
\begin{equation}\label{eq:expansion}
  w_{n,\,1/2+h}\;=\;\tfrac12+\alpha_{n}^{\bullet}h+\beta_{n}^{\bullet}h^{2}
  +\gamma_{n}^{\bullet}h^{3}+\delta_{n}^{\bullet}h^{4}+O(h^{5}),
  \qquad \bullet\in\{\text{left},\text{right}\},
\end{equation}
valid for $h\uparrow 0$ ($\bullet=\text{left}$) and $h\downarrow 0$
($\bullet=\text{right}$) respectively.

\section{Polynomials for small $n$}

At $p=\tfrac12$ every strategy is optimal, so the fully optimal value
agrees with strategy~B's value on a neighbourhood. Iterating
$b_{n}=p^{n}+\sum_{j=1}^{n-1}\binom{n}{j}p^{n-j}(1-p)^{j}b_{j}$ with
$b_{0}=1$ gives
\begin{align*}
  w_{1}(p) &= p,\\
  w_{2}(p) &= 3p^{2}-2p^{3},\\
  w_{3}(p) &= 13p^{3}-27p^{4}+21p^{5}-6p^{6},\\
  w_{4}(p) &= 75p^{4}-316p^{5}+606p^{6}-664p^{7}+432p^{8}-156p^{9}+24p^{10}.
\end{align*}
The degree of $w_n$ grows like $\binom{n+1}{2}$; these representations
match the optimal $w_{n}$ on a two-sided neighbourhood of $p=\tfrac12$
only for $n\le 2$, and on a left neighbourhood for those $n$ where
strategy~B remains optimal.

\section{Right-hand expansion ($p>\tfrac12$)}

By Manuscript Lemma~0.3, strategy~A is optimal for $p>\tfrac12$, so
$w_{n,p}=a_{n,p}$ where
\[
  a_{n}\;=\;p^{n}+\bigl(1-p^{n}-(1-p)^{n}\bigr)a_{n-1},\qquad a_{0}=1.
\]
Differentiating at $p=\tfrac12$ and using $p^{n}|_{1/2}=(1-p)^{n}|_{1/2}=1/2^{n}$ yields
\begin{equation}\label{eq:alpha_right}
  \alpha_{n}^{\text{right}}\;=\;\frac{n}{2^{n-1}}
  +\Bigl(1-\frac{1}{2^{n-1}}\Bigr)\alpha_{n-1}^{\text{right}},
  \qquad \alpha_{1}^{\text{right}}=1.
\end{equation}

\begin{proposition}
  $(\alpha_{n}^{\text{right}})_{n\ge1}$ is strictly increasing and
  converges to a finite limit $\alpha_{\infty}^{\text{right}}$.
\end{proposition}

\begin{proof}
  The forced positivity of $n/2^{n-1}$ and the contraction factor
  $1-1/2^{n-1}\in[0,1)$ force
  $\alpha_{n}^{\text{right}}-\alpha_{n-1}^{\text{right}}=
  n/2^{n-1}-\alpha_{n-1}^{\text{right}}/2^{n-1}$, which one checks is
  strictly positive by induction using
  $\alpha_{n}^{\text{right}}<n$ (immediate from \eqref{eq:alpha_right}).
  Boundedness ($\alpha_{n}^{\text{right}}\le n$ for all $n$ is loose;
  a sharper bound follows from Lemma~0.3 together with
  $w_{n,p}\le1$) gives convergence.
\end{proof}

Numerically,
\[
  \alpha_{1}^{\text{right}}=1,\quad
  \alpha_{2}^{\text{right}}=\tfrac{3}{2},\quad
  \alpha_{3}^{\text{right}}=\tfrac{15}{8},\quad
  \alpha_{4}^{\text{right}}=\tfrac{137}{64},\quad
  \alpha_{5}^{\text{right}}=\tfrac{2375}{1024},\quad\cdots\quad
  \alpha_{\infty}^{\text{right}}\approx 2.607.
\]

\section{Left-hand expansion ($p<\tfrac12$)}

For $h\uparrow 0$ the quantity $M_{n,j}$ in \eqref{eq:bellman} satisfies
\[
  M_{n,j}\bigl(\tfrac12+h\bigr)
  \;=\;\tfrac12\;+\;\min_{j\le m\le n-1}\alpha_{m}^{\text{left}}\cdot h\;+\;O(h^{2}),
\]
because among $w_{m}(\tfrac12+h)=\tfrac12+\alpha_{m}^{\text{left}}h+\cdots$ (for $h<0$),
the maximum is attained where $\alpha_{m}^{\text{left}}$ is
\emph{smallest}. Substituting into \eqref{eq:bellman} and using the
combinatorial identity
\[
  \sum_{j=1}^{n-1}\binom{n}{j}(n-2j)\;=\;0
\]
(obtained from $\sum_j\binom{n}{j}=2^{n}$ and
$\sum_j j\binom{n}{j}=n\,2^{n-1}$) yields the recursion
\begin{equation}\label{eq:alpha_left}
  \boxed{\;\alpha_{n}^{\text{left}}
  \;=\;\frac{n}{2^{n-1}}
  +\frac{1}{2^{n}}\sum_{j=1}^{n-1}\binom{n}{j}
  \min_{j\le m\le n-1}\alpha_{m}^{\text{left}},\qquad\alpha_{1}^{\text{left}}=1.\;}
\end{equation}

Numerical values (to 18 digits, via \eqref{eq:alpha_left} with mpmath):
\[
\begin{array}{c|l}
  n & \alpha_{n}^{\text{left}}\\\hline
  1 & 1.000000000000000000\\
  2 & 1.500000000000000000\\
  3 & 1.687500000000000000\\
  4 & 1.734375000000000000\\
  5 & 1.735839843750000000\\
  6 & 1.729385375976562500\\
  7 & 1.722763895988464355\\
  8 & 1.717139652930200100\\
  9 & 1.712851968293762180\\
  10& 1.709751579242984930\\
  \vdots & \vdots\\
  \infty & \approx 1.703
\end{array}
\]

\begin{proposition}\label{prop:alphaleft}
  The sequence $(\alpha_{n}^{\text{left}})_{n\ge1}$ is strictly
  increasing on $\{1,\dots,5\}$, strictly decreasing on $\{5,6,\dots\}$,
  and converges to a limit $\alpha_{\infty}^{\text{left}}\approx 1.703$.
  In particular it has a global maximum at $n=5$.
\end{proposition}

Numerical verification of Proposition~\ref{prop:alphaleft} up to $n=20$ is
straightforward and agrees with the table above.

\begin{corollary}[first-order shape of $n\mapsto w_{n,p}$ for
$p\lesssim\tfrac12$]
\label{cor:valley}
  For every sufficiently small $\varepsilon>0$ there exists $h_{0}<0$
  such that for all $h\in(h_{0},0)$,
  \[
    w_{1}\;>\;w_{2}\;>\;w_{3}\;>\;w_{4}\;>\;w_{5}\;<\;w_{6}\;<\;w_{7}\;<\cdots,
  \]
  up to corrections of order $|h|^{2}$. In particular $n\mapsto w_{n,p}$
  is valley-shaped at first order with a unique strict local minimum at
  $n=5$ and \emph{no} strict local maximum at first order.
\end{corollary}

\begin{proof}
  $w_{n,\,1/2+h}-w_{n-1,\,1/2+h}=
  (\alpha_{n}^{\text{left}}-\alpha_{n-1}^{\text{left}})\,h+O(h^{2})$.
  The sign of the difference is opposite to the sign of
  $\Delta\alpha_{n}:=\alpha_{n}^{\text{left}}-\alpha_{n-1}^{\text{left}}$.
  By Proposition~\ref{prop:alphaleft}, $\Delta\alpha_{n}>0$ for
  $n=2,\dots,5$ and $\Delta\alpha_{n}<0$ for $n\ge6$, which gives the
  claimed chain of inequalities for $h<0$ small.
\end{proof}

\begin{remark}
  Corollary~\ref{cor:valley} says that local maxima of $n\mapsto w_{n,p}$
  for $p<\tfrac12$ are \emph{not} visible at first order in $h=p-\tfrac12$.
  They arise only once $|h|$ is large enough that higher-order Taylor
  terms contribute; this matches the numerical observation that the first
  local maximum of $(w_{n,p})$ appears at increasingly large $n$ as
  $p\to\tfrac12^-$ (e.g.\ at $n=154$ for $p=0.49$, but at $n=9$ for
  $p=0.42$).
\end{remark}

\section{Higher-order coefficients (strategy~B)}

For $n$ up to the threshold where $w_{n,p}=b_{n,p}$ on a left
neighbourhood of $p=\tfrac12$, the Taylor expansion
\eqref{eq:expansion} has coefficients equal to those of the polynomial
$b_{n}(p)$, obtained by straight differentiation. Compactly, the first
four coefficients for $n=1,\dots,8$ are:

\[
\small
\begin{array}{c|cccc}
  n & \alpha_{n}^{\text{B}} & \beta_{n}^{\text{B}} & \gamma_{n}^{\text{B}} & \delta_{n}^{\text{B}}\\\hline
  1 & 1 & 0 & 0 & 0\\
  2 & 3/2 & 0 & -2 & 0\\
  3 & 27/16 & -3/8 & -7/2 & 3\\
  4 & 111/64 & -25/32 & -15/4 & 7\\
  5 & 3555/2048 & -535/512 & -1735/512 & 295/32\\
  6 & 113337/65536 & -19167/16384 & -3045/1024 & 39165/4096\\
  7 & 14471919/8388608 & -5067741/4194304 & -5624815/2097152 & 9399929/1048576\\
  8 & 462775983/268435456 & -161718519/134217728 & -168136367/67108864 & 274696555/33554432
\end{array}
\]
so that, for those $n$,
\[
  w_{n,\,1/2+h}\;=\;\tfrac12
  +\alpha_{n}^{\text{B}}h+\beta_{n}^{\text{B}}h^{2}
  +\gamma_{n}^{\text{B}}h^{3}+\delta_{n}^{\text{B}}h^{4}+O(h^{5}).
\]

\begin{remark}
  The first-order coefficients $\alpha_{n}^{\text{B}}$ coincide with
  $\alpha_{n}^{\text{left}}$ as long as the minimum in
  \eqref{eq:alpha_left} is attained at the \emph{left} endpoint
  ($m=j$) for every $j$, which is the case precisely when strategy~B is
  optimal at step~$n$. Once that fails (i.e.\ once the sequence
  $\alpha_{m}^{\text{left}}$ stops being monotone increasing), $\alpha^{\text{B}}$
  and $\alpha^{\text{left}}$ diverge. The first such discrepancy occurs
  between $n=5$ and $n=6$, which is exactly where the numerical
  monotonicity of the optimal $w$-sequence first breaks for $p$ close to
  $\tfrac12$.
\end{remark}

\section{Summary}

\begin{itemize}
  \item $w_{n,p}$ is $C^{0}$ but typically only piecewise $C^{\infty}$
  near $p=\tfrac12$; left/right derivatives differ for $n\ge3$.
  \item Right-side: $\alpha_{n}^{\text{right}}$ strictly increases,
  consistent with Lemma~0.3.
  \item Left-side: $\alpha_{n}^{\text{left}}$ is valley-shaped
  (increasing then decreasing), so at first order in $h=p-\tfrac12<0$ the
  sequence $(w_{n,p})_{n}$ has a strict local minimum at $n=5$ and no
  local maximum.
  \item Local maxima of $(w_{n,p})$ for $p<\tfrac12$ are therefore a
  \emph{non-perturbative} phenomenon near $p=\tfrac12$: they appear only
  at $n$ and $p$ at which higher-order Taylor terms overpower the
  first-order valley structure.
\end{itemize}

\end{document}
